\documentclass[11pt,a4paper]{article}
\usepackage[margin=2.2cm]{geometry}
\usepackage[T1]{fontenc}
\usepackage{amsmath,amssymb,xcolor,enumitem,booktabs,graphicx,tabularx}
\usepackage[normalem]{ulem}
\setlength{\parskip}{0.6em}\setlength{\parindent}{0pt}
\definecolor{newc}{HTML}{1f6f43}
\newcommand{\new}[1]{\textcolor{newc}{#1}}
\newcommand{\old}[1]{\textcolor{gray}{\sout{#1}}}
\newcommand{\note}[1]{{\small\itshape [#1]}}
\newcommand{\lab}[1]{\textsf{\footnotesize\detokenize{#1}}}
\renewcommand{\ref}[1]{\lab{#1}}
\newcommand{\citep}[1]{[\lab{#1}]}
\newcommand{\citet}[1]{\lab{#1}}
\title{Proposed change, Introduction and Project Description,\\ second pass}
\date{24 November 2026, version 1}
\begin{document}\sloppy
\maketitle

\section*{How to read this}
This answers the red comments on the Introduction and Project Description
(\texttt{intro\_comments.pdf}, 24 November). Agreed: the storage wording on
p.~1 must follow the arrival model of the code, and Token Passing needs a
reason for being the reference. Part~4 lists the values the chapter leans on,
where each is set and the reason the thesis gives, and asks where a reason is
missing.

Black text stays, \new{green} is new and \old{struck} text goes.
\note{Notes in brackets} say why. Nothing below is in the thesis yet.

\section*{Part 1. What was checked}
\begin{itemize}[nosep]
  \item \textbf{Storage, [A38] (2026-10-18).} The code changed from moving
    stock to placing arriving stock (\texttt{storage/arrival.py}: ``Nothing
    already on a shelf moves''). GAPS [A38] logs it, and
    \lab{sec:pf:scope} says ``Stored stock never moves''. The Background had
    the same leftover wording (``updates storage dynamically as demand
    shifts'') and lost it in [A48]. The Introduction still says ``move stock
    as demand changes''. It is a leftover, not a model change.
  \item \textbf{The three rules in the code.} Fixed storage sends a SKU's
    arriving units to its home shelf. Demand storage places the SKUs ordered
    most lately first, each on the free shelves nearest the workstations.
    Congestion-aware storage orders the shelves by a score that also counts
    how long picks stand blocked there.
  \item \textbf{Token Passing.} Appendix~E uses it for the load sweep
    ``because it needs no training, so the load does not depend on the
    training of a learned controller''. The Method section ``A reference
    point'' has a pending review (\texttt{review\_method\_reference.pdf} v2)
    that gives the same role. The sentence below matches it.
  \item \textbf{Results.} At the main load the learned controllers fall
    behind on most seeds. Only Token Passing keeps up on every seed.
\end{itemize}

% ================================================================
\section*{Part 2. Introduction}

An automated warehouse makes two kinds of decision that depend on each
other. \ldots The first is studied as the Storage Location Assignment
Problem (SLAP)\new{~\citep{reyes2019slap_review}}, including dynamic
versions where the placement changes over time. \ldots
\note{SLAP gets the same backing as MAPD. See question~1 on a sentence of
motivation before this paragraph.}

Storage and routing affect each other. \ldots So neither layer can be
planned well on its own.

This thesis measures how much it helps to \old{move stock as demand
changes}\new{place arriving stock by demand}, and whether that gain depends
on how the robots are controlled. \new{The gain is measured in service time,
the mean time from when an order is placed until it leaves the warehouse.}
A storage rule decides from time to time where arriving stock
goes\old{, based on the orders it has seen}. \new{Stock already on a shelf
stays there. Fixed storage sends every SKU to one home shelf. Demand storage
puts the SKUs ordered most lately on the free shelves nearest the
workstations.} Between these updates, a controller gives tasks to the
robots and routes them. Demand storage is compared with fixed storage under
three controllers. The centralised controller decides for the whole fleet,
the section-based one for one zone and the decentralised one for a single
robot. All three are graph neural networks trained the same way with
reinforcement learning, so they differ \old{only} in how much one decision
covers\new{. What each robot sees and whose reward judges its move follow
from that}. Token Passing, a classic planner that does not learn, is shown
next to them as a reference. \new{It comes from the paper that defines MAPD,
so it shows what a known planner reaches in the same warehouse. Because it
needs no training, it also sets the main load, the rate of orders at which
the comparison is made.}
\note{``move stock'' was a leftover from before [A38]. The two rules are
defined where they are first compared. The measure moves here from its own
one-line paragraph. ``only'' was not the whole story, since view and credit
differ too ([A128]), but both follow from the scope of the decision. The
Token Passing reason matches Appendix~E and the pending Method review, and
it defines ``main load'' before the next paragraph uses it.}

The coupling is tested in two ways. The first runs one way. \ldots If it
does, what storage is worth depends on how the robots are routed.

The second runs both ways. \new{A third rule, congestion-aware storage,}
\old{The storage rule} can also read where robots \old{get stuck}\new{are
blocked, standing still because the way ahead is taken}, and steer both new
stock and picks away from those shelves. Then the robots' movement feeds
back into where the stock goes and where it is taken from, and the loop
between the two layers is closed. Avoiding crowded aisles can only help
when the aisles are crowded. So the loop is tested at the main load and
again at a higher load.
\note{Congestion-aware storage is named before Aim and scope lists it.
``Blocked'' is the one name for the signal, also used by RQ2 and Expected
outcome.}

The decentralised robots may also send \old{short}\new{learned} messages to
the robots near them. A robot that decides alone sees only what is around
it, and a message can tell it about crowding further away.

\old{The main measure is the mean time an order takes, from when it is placed
until it leaves the warehouse.}
\note{Moved into the paragraph on what the thesis does.}

% ================================================================
\section*{Part 3. Project Description}

\subsection*{Aim and scope}

The project aims to find out how \old{feedback}\new{the coupling} between
storage and the control of the robots affects how \old{well a warehouse
works}\new{fast orders are served}.
\note{One name for the coupling, and the aim says what is measured.}

Four things are varied in the experiments: \ldots
\note{List unchanged.}

The main result compares demand storage with fixed storage under each
controller, at a main load where every controller should keep up. \new{The
main load is 0.6 of the highest rate at which Token Passing keeps up under
fixed storage. Near that limit a small saving in travel turns into a large
saving in waiting, and the gain would show how close the fleet runs to its
limit, not what storage adds.} The other factors test how far that result
reaches.
\note{The reason is the one Appendix~E gives. Where 0.6 comes from goes in
Appendix~E, see question~2.}

A fleet keeps up when it finishes tasks as fast as they arrive. Order times
are compared only when both runs keep up, because a fleet that falls behind
has a queue that grows for as long as the run lasts. \new{When a fleet falls
behind, it is judged by its served share, the part of the arriving work it
finishes.}
\note{RQ2 and RQ4 ask whether a fleet serves more. The served share was lost
when [A129] cut the measures paragraph.}

\subsection*{Research question}
\ldots
\begin{description}[nosep]
  \item[RQ3:] As orders arrive faster, does the gain of RQ1 grow, shrink or
  disappear, and at what rate does each \old{trained}\new{learned}
  controller stop keeping up?
\end{description}
\note{One name, ``learned'', as in the Introduction. Token Passing's limit
is already the capacity of Appendix~E, so leaving it out of RQ3 is right.}

\subsection*{Contributions}
\begin{enumerate}[nosep]
\item One model in which \old{the three problems}\new{storage assignment,
  task assignment and path finding} share the same warehouse graph and the
  same stock, so each changes what the others face.
\item Congestion-aware storage, the thesis's own rule. It adds the time
  \old{picks stand still}\new{robots stand blocked} at a shelf to the classic
  \old{slotting score}\new{score that puts the most ordered SKUs nearest the
  workstations~\citep{heskett1963coi}}, so where new stock goes reacts to
  the traffic of the \old{agents}\new{robots}.
\item A simulator built on that model. Between two runs \new{of one fleet
  size} only the storage rule or only the controller changes. \ldots
\end{enumerate}
\note{The three problems lost their names in [A129]. ``Slotting score'' was
used before the Background explains it. ``Robots'' in this chapter. The
warehouse grows with the fleet, so the pairing holds within one fleet size.}

\subsection*{Expected outcome}

\old{Before the experiments, the expectations are as follows.} For RQ1,
demand storage should shorten \old{order}\new{service} time under every
controller, mostly because popular stock sits closer to the workstations and
trips get shorter. The gain should be largest for the centralised
controller, because it sees where every \old{agent}\new{robot} is going
\ldots
\note{No sentence about the text. One name for the measure and for robots.
The RQ4 expectation is left as it is, see question~4.}

% ================================================================
\section*{Part 4. The values and their reasons}

\begin{tabularx}{\linewidth}{@{}p{2.6cm}p{2.2cm}X@{}}
\toprule
Value & Set in & Reason the thesis gives, and what is missing \\
\midrule
Token Passing as reference & App.~E, Method & Needs no training, so the
load does not depend on a learned controller. From the paper that defines
MAPD. Missing in the Introduction until now (Part~2). \\
Capacity & App.~E, [A58] & Highest tested rate at which all 20 seeds keep
up, and every lower rate too, so one odd seed is not read as the capacity.
Complete. \\
Main load, 0.6 of capacity & App.~E, [A58] & Kept below capacity, because
near the limit the gain would show the waiting, not storage. \textbf{Why 0.6}
is not in the thesis. Author, 24 November: it comes from an early Token
Passing run that found where the fleet falls behind ([A31] pilot, 2026-10-14,
3 agents, 6 seeds). The old appendix notes call 0.6 the point where the fleet
``keeps up, but agents don't wait for tasks''. Proposed text below. \\
Keep-up threshold, 0.9 & App.~E, [A111] & Token Passing serves 1.00 at the
main load and no seed under 0.94, so 0.9 does not fail by chance. 0.8 let
growing queues pass. Complete. \\
Higher load & [A84], TBD & Rule in [A84]: the highest load where every
\emph{trained} controller still keeps up. \textbf{It cannot be met.} The
learned controllers fall behind already at the main load. \\
\bottomrule
\end{tabularx}

% ================================================================
\section*{Questions}
\begin{enumerate}[nosep]
  \item Motivation. The Introduction opens on the two decisions and never
    says why warehouses or robot fleets matter. Add one or two sentences
    with a source (for example the survey of de Koster et al., 2007, which
    is in the bibliography), or leave it as it is?
  \item Why 0.6 (answered, please check the wording). Proposed for
    Appendix~E, after $\lambda_{\mathrm{main}} = 0.6\,\lambda_{\mathrm{capacity}}$:
    \new{The share 0.6 comes from an early Token Passing run on a small
    warehouse with three agents, made to find where the fleet falls behind.
    At about 0.6 of that point it still kept up with room to spare, and its
    agents seldom waited for a task. The share was kept when the warehouses
    grew, and the capacity of each fleet was measured again.} In Aim and
    scope only the clause ``so the fleet is busy but not at its limit''.
    Is ``seldom waited for a task'' true to what you saw? The [A31] entry
    does not record idle time, only where the fleet stopped keeping up.
  \item The higher load. The [A84] rule has no answer. Options:
    (a)~set it from Token Passing too, at a share between 0.6 and 1.0 of
    capacity, for example 0.9, and run all controllers there, judging the
    learned ones by served share; (b)~drop the higher load, since the
    learned controllers already meet crowded aisles at the main load, and
    RQ2 reads the main load only; (c)~keep it for Token Passing only, where
    a higher load still keeps up. I lean to (a), because RQ2 already asks
    what a fleet that falls behind serves. It changes [A84], the
    Introduction, RQ2 and the Experiments table, so it gets its own tag.
  \item RQ4 now asks whether the effect of messages changes with the storage
    rule, but Expected outcome predicts only that messages help. What did
    you expect for the storage rule before the experiments? I do not write
    an expectation after the result is known.
\end{enumerate}

\end{document}
